% Author: Keshav (condensed from report/section_parsing.md)
\section{Address parsing with a BiLSTM-CRF}
\label{sec:parsing}

Two branches of one chain often differ only in the house number or the street, so a matcher should be able
to compare addresses field by field. A parser labels each address token with one of the tags
\texttt{HOUSE}, \texttt{STREET}, \texttt{LOCALITY}, \texttt{LANDMARK}, \texttt{CITY}, \texttt{STATE}, \texttt{POSTCODE},
plus \texttt{SEP} for commas. This is sequence labelling, and a field is a maximal run of one tag:

\begin{center}\small
\begin{tabular}{cccccc}
\texttt{003rd floor} & \texttt{platinum plaza plot no.24 pu-4} & \texttt{behind c-21 mall} & \texttt{a.b.\ road} & \texttt{indore} & \texttt{mp} \\
\textsc{house} & \textsc{locality} & \textsc{landmark} & \textsc{street} & \textsc{city} & \textsc{state}
\end{tabular}
\end{center}

\paragraph{Silver labels.} There are no gold parses, so labels come from a rule parser that relies on the
commas between components. It finds the state at the end, the city before it, and leading house numbers,
and it tags landmark and street phrases by keywords. State spellings are learned on train: a component
that occurs at least 30 times and is the last component in at least 60\% of its occurrences. This gives 42
spellings, including romanised ones such as \emph{kerlm} for {\mlym കേരളം}. Without commas, the rules
cannot find component boundaries at all.

\paragraph{Model.} Each token is represented by a word embedding (100 dimensions; 30,000 most frequent
training tokens, digits mapped to 0) concatenated with a character BiLSTM ($2\times32$), which covers unseen
house numbers and spelling variants. A 2-layer BiLSTM ($2\times128$) produces emission scores for the 8 tags.
A CRF output layer scores whole tag sequences with learned transitions
(\textsc{city}$\to$\textsc{state} likely, \textsc{state}$\to$\textsc{house} rare), trained by maximum
likelihood and decoded with Viterbi. The model has 2.72M
parameters. \textbf{Training trick:} the silver labels are computed on the address with commas. Then, for
half of the training addresses, the comma tokens are deleted from both tokens and tags, so the model must
learn the boundaries from the words themselves. It is trained on 120,000 training addresses for 3 epochs
(Adam, 1.4 minutes on a T4).

\begin{table}[h]
\centering\small
\begin{tabular}{llccc}
\toprule
input & parser & token accuracy & macro F1 & exact address \\
\midrule
with commas    & rules      & 1.0000 & 1.0000 & 1.0000 \\
with commas    & BiLSTM-CRF & 0.9994 & 0.9989 & 0.9969 \\
commas removed & rules      & 0.5758 & 0.3691 & 0.0000 \\
commas removed & BiLSTM-CRF & \textbf{0.9500} & \textbf{0.9063} & \textbf{0.8545} \\
\bottomrule
\end{tabular}
\caption{Parsing 10,000 val addresses, against the silver labels. Macro F1 over the six tags that occur.}
\label{tab:parser}
\end{table}

\paragraph{Results.} With commas, the rules score 1.0 by construction and the BiLSTM-CRF reproduces them on
99.7\% of addresses (Table~\ref{tab:parser}). The real test is the address without commas. There the rules
put everything after the house number into one field and get no address fully right, while the BiLSTM-CRF
parses 85\% of addresses exactly. \textsc{state} (F1 0.999), \textsc{house} (0.979) and \textsc{city} (0.972)
are easy. \textsc{street} vs.\ \textsc{locality} is the hard boundary, and \textsc{landmark} has low precision
(F1 0.639), partly because the model finds landmark phrases in the middle of a component that the rules miss.

\paragraph{Limitations.} The parser is evaluated against the rules that trained it, so the numbers measure
how well it recovers the comma structure, and rule errors propagate. None of the 615,349 records has a
postcode the rules recognise, so \textsc{postcode} has no support. Every non-empty address in our data has
commas, so the comma-free condition is a stress test we created, not a case in this data.
\S\ref{sec:classical} measures whether the parser helps matching.
